\documentclass[12pt]{article}
\usepackage[utf8]{inputenc}
\usepackage[T1]{fontenc}
\usepackage{hyperref}
\title{Southern Rincon Zapotec Language Resources for K-12 Curriculum Development}
\author{Compiled by Claude Sonnet 5.5 (claude-sonnet-5-5) for Andr\'as Kornai}
\date{}
\begin{document}
\maketitle

\section{Introduction}
Southern Rincon Zapotec (ISO 639-3 zsr, Glottocode sout3004) is an Otomanguean language of Oaxaca, Mexico. ASJP gives about 12,000 speakers. Glottolog marks its status as ``shifting''. It is described as easily intelligible with Rinc\'on Zapotec (zar), but is classed as a separate language. Few resources turned up, and nothing for schools.

\section{Primary Educational Resources}
None found in the search results.

\section{Linguistic and Reference Materials}
\begin{itemize}
\item Glottolog entry for Southern Rincon Zapotec. \url{https://glottolog.org/resource/languoid/iso/zsr}
\item ASJP wordlist and data page. \url{https://asjp.clld.org/languages/ZAPOTEC_SOUTHERN_RINCON}
\item Wikitongues language page. \url{https://wikitongues.org/languages/zsr/}
\end{itemize}

\section{Religious and Community Resources}
\begin{itemize}
\item New Testament translation (1992, Wycliffe Bible Translators), available as a PDF and reported to be under a Creative Commons license. \url{https://nobanners.classicistranieri.com/biblepdf/zsrNT/zsrNT_all.pdf}
\end{itemize}

\section{Digital and Audio Resources}
\begin{itemize}
\item Global Recordings Network page for the language. The page indicates that no recordings are available. \url{https://www.globalrecordings.net/en/language/zsr}
\end{itemize}

\section{Curriculum Development Context}
No evidence of school use was found. The New Testament is the only substantial text found and could serve as a corpus. Any curriculum would need to be developed with community speakers. Materials for the closely related Rinc\'on Zapotec may be worth checking, but their suitability would have to be confirmed.

\end{document}